\documentclass[12pt]{article}
\usepackage[utf8]{inputenc}
\usepackage[T1]{fontenc}
\usepackage[margin=1in]{geometry}
\usepackage{setspace}
\usepackage{xcolor}
\usepackage[colorlinks=true, urlcolor=blue!60!black, linkcolor=blue!60!black]{hyperref}
\setlength{\parindent}{0em}
\setlength{\parskip}{0.45em}
\renewcommand{\baselinestretch}{1.0}
\begin{document}
\begin{center}
{\Large \textit{The Slaughter-House Cases}, 83 U.S. 36 (1873)\par}
\vspace{0.3em}
{\normalsize Case Brief --- One-Page Edition\par}
\end{center}

\textbf{Memory Jogger.} The case that gutted the Privileges or Immunities Clause while upholding a New Orleans slaughterhouse monopoly (\href{https://www.courtlistener.com/opinion/88661/butchers-benevolent-assn-v-crescent-city-live-stock-landing/}{83 U.S. 36}).

\textbf{Facts.} For public health, Louisiana gave one company a 25-year exclusive right to slaughter livestock around New Orleans, requiring it to admit all butchers at capped fees. Displaced butchers sued under the Thirteenth and Fourteenth Amendments.

\textbf{Procedural History.} The Louisiana Supreme Court upheld the law; the butchers brought writs of error.

\textbf{Votes.} Affirmed, 5--4: Miller for the Court; Field, Bradley, and Swayne dissented, with Chase.

\textbf{Holding.} The Clause protects only rights of national citizenship, not state-citizenship rights such as pursuing a trade; the monopoly was neither servitude nor a due process or equal protection violation.

\textbf{Reasoning.} Section 1 distinguishes national from state citizenship; a broader reading would make the Court a ``perpetual censor upon all legislation of the States.'' Field said the majority made the Clause ``a vain and idle enactment, which accomplished nothing''; Bradley protected the right to choose a calling.

\textbf{Examples.} \emph{Same-Side:} a single waste-hauling franchise. \emph{Opposite-Side:} a tax on travel to the capital. \emph{Fence-Sitter:} incorporation through the Clause.

\textbf{Critique.} Originalists see a textual error that erased the Clause; progressives defend the health-law result but lament the retreat from Reconstruction.

\textbf{Key Quotation.} ``[N]o one can fail to be impressed with the one pervading purpose found in them all, lying at the foundation of each, and without which none of them would have been even suggested; we mean the freedom of the slave race, the security and firm establishment of that freedom, and the protection of the newly-made freeman and citizen from the oppressions of those who had formerly exercised unlimited dominion over him.'' --- Justice Miller, for the Court, 83 U.S. at 71.

\end{document}
